% Generated draft: requires mathematical and source review.
\documentclass[11pt]{article}
\usepackage[margin=25mm]{geometry}
\usepackage{fontspec}
\setmainfont{Latin Modern Roman}
\usepackage{amsmath,amssymb,mathtools,unicode-math}
\setmathfont{Latin Modern Math}
\usepackage{enumitem,needspace,xurl,hyperref}
\setlength{\parindent}{0pt}
\setlength{\parskip}{5pt}
\setlength{\emergencystretch}{3em}
\newcommand{\N}{\mathbb{N}}
\newcommand{\Z}{\mathbb{Z}}
\newcommand{\Q}{\mathbb{Q}}
\newcommand{\R}{\mathbb{R}}
\newcommand{\C}{\mathbb{C}}
\newcommand{\F}{\mathbb{F}}
\newcommand{\A}{\mathbb{A}}
\newcommand{\PP}{\mathbb{P}}
\newcommand{\E}{\mathbb{E}}
\newcommand{\eps}{\varepsilon}
\newcommand{\dd}{\,\mathrm{d}}
\begin{document}
\section*{UnsolvedMath ID: 30005764}
\textbf{Source problem code:} catalogue entry
\subsection{Definitions and mathematical statement}
A permutation of size \(n\) is a bijection \(\sigma:\{1,\ldots,n\}\to\{1,\ldots,n\}\). A pattern of \(\sigma\) is the permutation obtained by replacing the entries of a subsequence by their relative ranks. A permutation class \(C\) is a set of permutations closed under taking patterns. Write \(C_n=C\cap S_n\), and, whenever \(C_n\ne\varnothing\), let \(\sigma_{n,C}\) be uniformly distributed on the finite set \(C_n\). An asymptotic question as \(n\to\infty\) requires \(C\) to contain permutations of arbitrarily large size; a finite class does not provide such a sequence.

A permuton is a Borel probability measure on \([0,1]^2\) whose two coordinate marginals are uniform. Associate to \(\sigma\in S_n\) the permuton \(\mu_\sigma\) with density \(n\) on the union of the squares
\[
\bigcup_{i=1}^n
\left[\frac{i-1}{n},\frac{i}{n}\right]
\times
\left[\frac{\sigma(i)-1}{n},\frac{\sigma(i)}{n}\right],
\]
and density zero elsewhere, with square boundaries irrelevant to the measure. Equip the space of permutons with the weak topology.

\textbf{Permuton convergence question.} For every permutation class \(C\) containing permutations of arbitrarily large size, does \(\mu_{\sigma_{n,C}}\) converge as \(n\to\infty\)? The supplied record does not specify whether the intended conclusion is convergence in distribution to a possibly random permuton,
\[
\mu_{\sigma_{n,C}}\xrightarrow{\mathrm d}M_C,
\]
or convergence in probability to a deterministic permuton \(\mu_C\). These are distinct formulations; neither is selected here. Almost-sure convergence would additionally require a coupling across sizes, which is not supplied.

\textbf{Graphon direction.} A graphon, in the usual simple undirected graph setting, is a symmetric measurable function \(W:[0,1]^2\to[0,1]\); graph limits are commonly considered up to zero cut distance. The source asks what happens in contexts such as graphons. It does not specify the graph families, their closure properties, the sampling law, the limiting topology, or whether the limit is deterministic or random. Thus this is an open direction in the supplied record, rather than a fully quantified graphon conjecture. The preceding definition describes a conventional setting and does not establish that this is the intended setting.

\textbf{Algebraicity questions.} The coefficient field, its characteristic, and the admissible formal-series rings are not specified. To express the requested implication, let \(K\) denote a field that must be specified as part of a precise formulation. In the possible formal interpretation \(F\in K[[t,u]]\), the series \(F(t,0)\) and \((\partial_uF)(t,0)\) are well defined in \(K[[t]]\). Here algebraicity over \(K(t,u)\) means that there exists a nonzero polynomial \(Q\in K[t,u,Y]\) such that
\[
Q(t,u,F(t,u))=0.
\]
The first requested task is to identify explicit hypotheses on the coefficient setting, a polynomial \(P\), and its admissible solutions under which
\[
P\bigl(F(t,u),F(t,0),(\partial_uF)(t,0),t,u\bigr)=0
\quad\Longrightarrow\quad
F(t,u)\text{ is algebraic over }K(t,u).
\]
For this formal interpretation one may take \(P\in K[X,A,B,t,u]\). This ambient ring is an explicit interpretation, not a coefficient specification supplied by the source. The original excerpt includes ``etc.'' after the listed arguments, so the displayed equation need not exhaust the intended types of auxiliary terms. No unconditional implication for arbitrary \(P\) or arbitrary solutions is asserted.

For the multivariate question, a corresponding possible interpretation uses finite lists
\[
F\in K[[t,u,v]],\qquad
G_1,\ldots,G_r\in K[[t,u]],\qquad
H_1,\ldots,H_s\in K[[t]],
\]
and a polynomial relation
\[
P\bigl(F(t,u,v),G_1(t,u),\ldots,G_r(t,u),
H_1(t),\ldots,H_s(t),t,u,v\bigr)=0.
\]
Identify explicit hypotheses under which every admissible solution \(F\) is algebraic over \(K(t,u,v)\), meaning that some nonzero \(Q\in K[t,u,v,Y]\) satisfies \(Q(t,u,v,F)=0\). The source does not specify relations determining the auxiliary series, initial conditions, uniqueness requirements, or nondegeneracy assumptions. Its phrase ``natural assumptions'' remains an unresolved part of the requested characterization, not an established hypothesis or theorem.

\textbf{Extraction boundary.} The original excerpt also contains questions about uniquely-sorted permutations and pattern avoidance, followed later by a separately headed shuffle conjecture attributed there to Mathilde Bouvel. Their relationship to this registry entry is unclear. They are retained as source-context issues in the accompanying notes and are not silently incorporated into the two targets above.
\subsection{Short English statement}
The supplied entry asks whether uniform permutations from every class of unbounded size have a permuton limit, seeks an analogous graphon question, and asks which polynomial functional equations guarantee algebraicity of their formal-series solutions. The convergence mode, graphon model, coefficient setting, and algebraicity hypotheses remain unspecified. This is a statement-only draft requiring human mathematical and source review; the supplied open-status assessment is dated curation.
\subsection{Sources}
\begin{itemize}
\item Supplied local UnsolvedMath registry record, permanent ID \(30005764\), source problem code \texttt{catalogue entry}, title ``Universal Permuton Convergence and Algebraic Functional Equations,'' in the Oberwolfach Reports Open Problems collection. The record gives proposed year \(2024\) and no registry-level proposer.
\item \emph{Mini-Workshop: Permutation Patterns} (2024), Oberwolfach Reports, pp.\ 273--308, DOI \(10.4171/owr/2024/6\): \url{https://doi.org/10.4171/owr/2024/6}. This is also the supplied literature-source link labelled ``original OWR report.'' The supplied background describes these reports as documenting mathematical workshops and their open-problem discussions. The extracted text places the heading ``7. Mireille Bousquet-M\'elou'' immediately before the convergence and functional-equation discussion; the precise attribution boundaries require review.
\item Supplied literature source labelled ``graphon--permuton paper'': \url{https://arxiv.org/abs/1307.2444}. No author, title, or specific supporting result for this reference is supplied.
\item Literature assessment date: \texttt{2026-08-22}. This is supplied dated curation, not independently verified here. The record labels the problem open and reports that the graphon extension is posed as an open direction, that related graphon and permuton theory exists, and that no general answer was located. The supplied assessment lists the preceding two links as checked and expressly calls for independent verification of time-sensitive status and cited claims.
\item The registry labels its statement \texttt{corrected\_verified} and reports that the official report was checked directly and that the cleaned statement restores definitions and removes extraction artifacts. These are supplied verification claims, not checks performed for this draft.
\end{itemize}
\end{document}
